% Incorporación TER-01--04, TER-06--07. TER-05 conserva su desarrollo en
% 09_area, 09b_hojas_memoria, 09c_entropia_sectorial y la composición posterior.
% Fuentes: BOLTZMANN_DESARROLLO_OPERATORIO_20260927/{retorno,normalizacion,
% incidencia,seleccion,accion_horizonte_20260927}; FOCK_PARIDAD_RAMAS_Y_SELECTOR;
% REVISION_LIFT_POSITIVO_Y_ESPECTRO; ADDENDUM_TRANSPORTE_TERMICO_Y_SECCIONES.
% Edición aditiva: no sustituye el núcleo ni las demostraciones anteriores.
% Recibo causal heredado: APP conserva suma, producto, residuos y cocientes.
% TRIT conserva régimen y orientación en el calendario recibido.
% TPK transporta ruta, carry, frontera y memoria en la conexión nonádica.
% El estado enriquecido prolonga el registro al retornar la fase.
% La estructura discreta conjunta del continuo es el dominio heredado.
% Salida HMT examinada: capacidad, portador graduado y lector marcado.
% La realización espectral convencional es posterior y conserva las reglas declaradas.
\section{Capacidad espectral, memoria energética y transporte térmico}
\label{sec:ii-ter27-antecedentes}

La construcción siguiente reúne la capacidad de las prolongaciones, la
renovación de las dos hojas y sus realizaciones espectrales. Su objeto es
determinar qué conserva una lectura térmica cuando la fase retorna y el
registro sigue aumentando. Las operaciones proceden del estado enriquecido
APP--TRIT--TPK: el catálogo de firmas, el calendario, la orientación y la
memoria anteceden a las trazas que se evaluarán. Los espacios lineales,
las compresiones y las preparaciones que se añaden a esos antecedentes
se definen expresamente. Cada resultado conserva así sus condiciones de
construcción y su función en el argumento.

\subsection{Capacidad, vacancia y renovación conjunta}
\label{sec:ii-ter27-capacidad}

La comparación entre las $729$ palabras de seis trits y las $1000$
posiciones de un bloque decimal define, con el origen ya fijado en
\ref{vac-capacidad-trit},
\begin{equation}
 C_t=\lfloor(t+1)\rho\rfloor,
 \qquad \rho=\log_{1000}729=\log_{10}9,
 \qquad z_t=1-(C_t-C_{t-1}).
 \label{eq:ii-ter27-capacidad}
\end{equation}
La desigualdad $0<\rho<1$ implica $z_t\in\{0,1\}$. La capacidad cuenta
bloques en esta comparación; su evaluación exacta se efectúa mediante
$10^{C_t}\leq9^{t+1}<10^{C_t+1}$. Para $1\leq t\leq20$ se obtiene
$C_t=t$, mientras $C_{21}=20$. En efecto, para $t\leq20$ la desigualdad
inferior equivale a $9(9/10)^t\geq1$, y basta comprobarla en $t=20$;
en $t=21$ esa desigualdad se invierte. Las dos comparaciones restantes
sitúan $9^{22}$ entre $10^{20}$ y $10^{21}$. La primera vacancia positiva
es, por tanto, $p_1=21$, con
\[
 C_{20}=C_{21}=20,\qquad C_{22}=21,\quad C_{23}=22,\quad C_{24}=23.
\]
Los veinte incrementos de capacidad forman veinte etiquetas efectivas.
El índice de corte $24$, las $25$ posiciones de $0$ a $24$ y la capacidad
$23$ son registros de tipos distintos.

El catálogo conjunto de emisiones es
$\Omega=\mathcal S_+\times\mathcal S_\times$, con cardinales $18$ y
$26$. Sobre $\ell^2(\Omega)$ con conteo uniforme, las renovaciones de
la coordenada activa son
\[
 E_+=\frac{\mathbf J_{18}}{18}\otimes I_{26},\qquad
 E_\times=I_{18}\otimes\frac{\mathbf J_{26}}{26}.
\]
Son proyectores ortogonales que conmutan. Su producto es el proyector
constante $\Pi$, cuya matriz tiene todas las entradas iguales a $1/468$.
La lectura reducida de fase utiliza el régimen de origen:
\[
 K_{\rm ph}((u,r),(v,r+1))=P_{\tau(r)}(u,v),
 \qquad (P_+,P_-,P_0)=(E_+,E_\times,I),
 \quad \tau(r)=M_{\rm ph}(1+r).
\]
Esta operación es posterior a la dinámica enriquecida y a su holonomía,
en el orden $U_t\longrightarrow\Gamma_9\longrightarrow K_{\rm ph}$.

\begin{proposition}[Renovación de las dos hojas sobre el reloj de capacidad]
\label{prop:ii-ter27-renovacion}
Evalúese el lector anterior mediante $\mathcal R_t=K_{\rm ph}^{a_t}$,
donde $a_t=C_{t+1}-C_t$. Desde la primera vacancia, el primer instante en
que su salida matricial es independiente de la emisión inicial es $24$.
Los rangos sucesivos son $468,26,1$; en la orientación conjugada son
$468,18,1$.
\end{proposition}
\begin{proof}
Las capacidades de origen $20,21,22$ corresponden a los regímenes
$0,+1,-1$. Los avances $21\to22\to23\to24$ aplican, por tanto,
$I,E_+,E_\times$, y sus productos acumulados son $I,E_+,\Pi$.
El rango de un producto tensorial es el producto de sus rangos:
$\operatorname{rank}E_+=26$, $\operatorname{rank}E_\times=18$ y
$\operatorname{rank}\Pi=1$. Las filas de $I$ y $E_+$ dependen de la
emisión inicial; las de $\Pi$ son idénticas. Esto prueba la minimalidad.
La secuencia conjugada aplica $I,E_\times,E_+$ y tiene el mismo producto
final. La composición por capacidad conserva el tiempo cronológico:
\[
 (t,C_t,v_t)\mapsto(t+1,C_t+a_t,v_t+1-a_t),\qquad v_t=t-C_t.
\]
La identidad $t=C_t+v_t$ se mantiene también cuando $a_t=0$.
\end{proof}

El descenso de la fase módulo nueve al régimen módulo tres es exacto
para este lector. Para $Q(u,r)=(u,[r]_3)$,
\[
 \sum_{s:[s]_3=b}K_{\rm ph}((u,r),(v,s))
 =\mathbf1_{b=[r+1]_3}P_{\tau(r)}(u,v).
\]
La suma depende sólo de $Q(u,r)$, lo que prueba la agregación fuerte.
En cambio, el levantamiento conserva $C=9m+3b+a$. El paso $20\to23$
transforma $(m,b,a)=(2,0,2)$ en $(2,1,2)$: retorna $a$ y avanza $b$.
La fase positiva módulo nueve pasa de $3$ a $6$. La renovación uniforme
es una salida matricial de rango uno; la ruta y el microestado previo se
conservan en el archivo enriquecido, no en esa salida matricial sola.
La información de un régimen uniforme es $\log3$; la de las emisiones
renovadas es $\log468$; la tasa de una rotación determinista es cero.
Estas tres lecturas corresponden a ensambles y operaciones diferentes.

\subsection{Retornos por capacidad y valoración de los cilindros}
\label{sec:ii-ter27-cociclo}

Defínase el retorno estricto del régimen por
\[
 R_3(s)=\min\{t>s:C_t>C_s,\ C_t\equiv C_s\pmod3\}.
\]
La sucesión $C_t$ es no acotada y sus incrementos son cero o uno.
Por ello visita $C_s+3$ antes que cualquier otro valor positivo congruente
con $C_s$; así $C_{R_3(s)}=C_s+3$. Desde $s_0=24$ se obtiene la
localización corregida
\begin{equation}
 s_j=\min\{t\geq24:C_t=23+3j\},\qquad C_{s_j}=23+3j.
 \label{eq:ii-ter27-localizacion}
\end{equation}
Los primeros tiempos son $24,27,30$. Los siguientes intervalos pueden
contener vacancias; el aumento de capacidad sigue siendo exactamente tres.

Para explicar la valoración, descompóngase cada bloque como
$d_r=100a_{r,1}+10a_{r,2}+a_{r,3}$, con $0\leq a_{r,k}\leq9$.
Un prefijo de $c$ bloques se transporta a los tres enteros
$X_k=\sum_{r=1}^c a_{r,k}10^{c-r}$. La división euclídea proporciona
una biyección; su inversa reagrupa las cifras de cada posición. El
truncamiento de un bloque se convierte en el truncamiento de una cifra
en cada coordenada. La suma con acarreo se transporta mediante esta
biyección completa, conservando sus cocientes.

El cilindro escalar de longitud $1000^{-c}$ corresponde al cubo semiabierto
\[
 \prod_{k=1}^3[X_k10^{-c},(X_k+1)10^{-c}),\qquad
 J_c=10^{-c}I_3,\qquad \det J_c=1000^{-c}.
\]
La igualdad expresa una correspondencia de medidas de cilindros. Entre
dos capacidades, $J_{c',c}=10^{-(c'-c)}I_3$ satisface
$J_{c'',c'}J_{c',c}=J_{c'',c}$. La similitud positiva isotrópica queda
determinada por su determinante, pues $s^3=1000^{-(c'-c)}$ tiene una
única raíz positiva. En los retornos marcados,
\begin{equation}
 J_{C_{s_j}}=10^{-23}1000^{-j}I_3.
 \label{eq:ii-ter27-cociclo}
\end{equation}
El factor $10^{-23}$ es la escala de una coordenada del cubo; la medida
del cilindro escalar a esa capacidad es $10^{-69}$. La distinción conserva
el grado geométrico del lector. Un extractor fijado en $s_0$, seguido de
truncamiento compatible, satisface $E_n\pi_{m,n}=E_m$ para $m\geq n\geq
s_0$; la prolongación conserva así el registro marcado sin recalcular sus
rangos a partir de un corte posterior.

\subsection{Portadores de relaciones y realización mediante intercambio}
\label{sec:ii-ter27-portador}

Sean $U_m$ el espacio de etiquetas de capacidad y $P_p$ el de posiciones
del prefijo. La expectativa diagonal sobre $\operatorname{End}(U_m)$ es
$E_{\rm diag}(X)=\sum_iE_{ii}XE_{ii}$. Su complemento ortogonal tiene
base $E_{ij}$, $i\ne j$, y dimensión $m(m-1)$. El segundo espacio de
registros es $\operatorname{End}(P_p)\oplus A_{p-1}$, donde $A_{p-1}$
conserva las transiciones adyacentes como observaciones tipadas. Su
dimensión es $p^2+p-1$. Con $m=20$, $p=25$, los rangos son $380$ y $649$.
Esta elección de espacios y de corte define una realización del lector;
las identidades siguientes demuestran sus operaciones sin atribuir al
conteo, por sí solo, la selección física de ese dominio.

La reversión basada actúa por
$E_{ij}\mapsto E_{m-1-i,m-1-j}$ y, en el segundo espacio, por
$E_{ij}\mapsto E_{p-1-i,p-1-j}$, $a_j\mapsto a_{p-2-j}$.
Para $m=20$ hay $190$ órbitas de dos elementos. Para $p=25$ existe un
único elemento matricial fijo, $E_{12,12}$; quedan $312$ órbitas
matriciales y $12$ órbitas de transiciones. Se obtiene una línea fija y
$324$ representaciones regulares de $C_2$, con autoespacios de dimensiones
$325$ y $324$. Éste es el antecedente de $649=1+2\cdot324$; conserva la
involución de registros, no determina una graduación energética.

Una segunda realización utiliza el mismo intercambio $J$ de dos
orientaciones y conserva los términos tensoriales cruzados. Sean $W,V$
espacios de dimensiones $m,n$. En el sector exterior de $W_+\oplus W_-$
y en los osciladores simétricos de $V_+\oplus V_-$, los caracteres fijos son
\begin{align}
 \chi_F(q)&=\tfrac12\bigl((1+q)^{2m}+(1-q^2)^m\bigr),\\
 \chi_B(q)&=\tfrac12\bigl(P_n(q)^2+P_n(q^2)\bigr),
 &P_n(q)&=\prod_{r\geq1}(1-q^r)^{-n}.
 \label{eq:ii-ter27-caracteres}
\end{align}
\begin{proof}
En una partícula, los autoespacios $J=\pm1$ tienen ambos dimensión $m$.
La traza insertada en el álgebra exterior es
$(1+q)^m(1-q)^m=(1-q^2)^m$. El proyector $(I+J)/2$ promedia esa traza
y el carácter ordinario. El coeficiente de grado dos es
$\bigl(\binom{2m}{2}-m\bigr)/2=m(m-1)$.
Para cada frecuencia bosónica, las multiplicidades de ambos signos son
$n$, por lo que la traza insertada es
$\prod_r[(1-q^r)(1+q^r)]^{-n}=P_n(q^2)$.
Los coeficientes ordinarios en grados cero, uno y dos son
$1,2n,2n^2+3n$, y los insertados $1,0,n$. Sus promedios son
$1,n,n(n+2)$. Cada coeficiente formal depende de un número finito de
factores, lo que justifica el producto sin hipótesis analítica adicional.
\end{proof}
El término adicional $n$ en energía dos procede de las creaciones de
frecuencia dos. En efecto,
$M(V)_2=V(-2)\oplus\operatorname{Sym}^2(V(-1))$; para rango $24$,
su dimensión es $24+300=324$. En las dos orientaciones intercambiadas,
el fijo bosónico acumulado hasta energía dos tiene dimensión
$1+n+n(n+2)=(n+1)^2+n$, que da $649$ para $n=24$.
La identificación lineal con el espacio previo de matrices y transiciones
requiere transportar sus bases; la igualdad de dimensiones no identifica
automáticamente sus Hamiltonianos. El corte exterior homogéneo y el corte
bosónico acumulado se mantienen declarados como partes de este lector.

\subsection{Semillas multiplicativas y compensación exacta de fibras}
\label{sec:ii-ter27-fibras}

El soporte de semillas multiplicativas es
$X_\times=(\mathbb Z/9\mathbb Z)^2\times\{N,E,S,O\}$, de cardinal
$324$. El emisor de seis ventanas $e:X_\times\to\mathcal S_\times$
del núcleo produce $26$ firmas. Sus fibras tienen cardinales $8$ para
$24$ firmas, $24$ para $555555$ y $108$ para $000000$; la suma es $324$.
El conjunto temporal $I=\{1,\ldots,24\}$ conserva un tipo distinto de la
fibra de $555555$, aunque ambos tengan cardinal $24$.

Se construye el espacio microscópico
\[
 \widetilde T=\{\Omega\}\sqcup\{a_i:i\in I\}
 \sqcup\{b_{i,y}:i\in I,y\in X_\times\},\qquad |\widetilde T|=7801.
\]
La lectura $p(b_{i,y})=b_{i,e(y)}$, que fija los otros registros, tiene
imagen $T$ de cardinal $1+24+24\cdot26=649$. Sea $m_e=|e^{-1}(e)|$.
Asignar masa $1$ a $\Omega,a_i$ y $1/m_{e(y)}$ a $b_{i,y}$ define
una medida $\nu$ cuyas fibras por $p$ tienen todas masa uno. Para
$r>0$, las tasas son $r$ en ambas direcciones entre $\Omega$ y $a_i$,
y $r/m_{e(y)}$ de $a_i$ a $b_{i,y}$, con tasa inversa $r$.

\begin{proposition}[Transporte compensado del árbol de firmas]
\label{prop:ii-ter27-arbol}
El generador microscópico es reversible para $\nu$. La aplicación $p$
induce el árbol simétrico sobre $T$ y el levantamiento $J_pg=g\circ p$
es una isometría que entrelaza ambos generadores.
\end{proposition}
\begin{proof}
En cada arista terminal se cumple el balance
$\nu(a_i)r/m_{e(y)}=\nu(b_{i,y})r$. Además,
$\sum_{y:e(y)=e}r/m_e=r$, y la tasa inversa tiene el mismo valor desde
cualquier representante. Esto prueba la agregación fuerte. La identidad
$\sum_{x:p(x)=z}\nu(x)=1$ prueba $J_p^*J_p=I$; las sumas de tasas
prueban $\widetilde L J_p=J_pL_T$. El árbol es conectado y tiene
$648$ aristas, de modo que su medida estacionaria normalizada es uniforme.
\end{proof}
La media vuelta del cursor induce una involución de firmas que conserva
$m_e$. Combinada con $i\mapsto25-i$, conserva $\nu$ y las tasas; el
entrelazador transporta también esa inversión. Las semillas permanecen
registradas en $\widetilde T$ aunque $T$ publique sólo su firma.
La selección de una marca y el refresco compensado son las operaciones
de este lector. Su dominio y sus tasas permanecen explícitos: la ruta
cardinal de un cursor y la actualización enriquecida completa conservan
operaciones adicionales. El árbol relaciona los niveles internos del
sector bosónico; las comunicaciones entre los tres grados se construirán
en \ref{sec:ii-ter27-corrientes}.

\subsection{Doble graduación y recuperación de sus marginales}
\label{sec:ii-ter27-doble}

El portador seleccionado para la comparación es
\[
 \mathcal D=\mathbb C\Omega_{\rm ext}\oplus F_2^J\oplus B_{\leq2}^J.
\]
El operador $L$ conserva el grado de Fock; $N$ asigna grados externos
$0,1,2$ a los tres sumandos. Con $m=20,n=24$, su descomposición conjunta es
\begin{equation}
 (N,L;d)=(0,0;1),(1,2;380),(2,0;1),(2,1;24),(2,2;624).
 \label{eq:ii-ter27-celdas}
\end{equation}
Los dos vacíos son ortogonales y conservan su procedencia. La función
bivariada y sus especializaciones son
\begin{align}
 \mathcal Z(u,v)&=1+380uv^2+u^2(1+24v+624v^2),\\
 Z_N(q)&=\mathcal Z(q,1)=1+380q+649q^2,\\
 Z_L(q)&=\mathcal Z(1,q)=2+24q+1004q^2.
 \label{eq:ii-ter27-particiones}
\end{align}
Cada monomio resulta de aplicar $u^Nv^L$ a una celda de
\eqref{eq:ii-ter27-celdas} y tomar su traza. En $q=1/1000$ se obtienen,
respectivamente, $1.380649$ y $2.025004$. La segunda evaluación reúne dos
vacíos, $24$ estados de energía uno y $1004$ de energía dos.

Las marginales completas $n_i,l_j$ de una distribución soportada en
esas cinco celdas recuperan sus masas:
\[
 p_{00}=n_0,\quad p_{12}=n_1,\quad p_{21}=l_1,
 \quad p_{20}=l_0-n_0,\quad p_{22}=l_2-n_1.
\]
La suma por filas y columnas prueba la restitución siempre que las
diferencias sean no negativas y los totales coincidan. El soporte es un
árbol bipartito; la restitución se refiere a sus cinco masas, mientras
las coherencias y los índices internos permanecen en el estado completo.
Dos cifras de partición aisladas no sustituyen esas marginales. Como
control, $Z_N-Z_L=(q-1)(1-355q)$: en $q=1/355$ ambas trazas coinciden,
pero las probabilidades del vacío interno son $1/261574$ y
$126025/261574$. La igualdad del escalar conserva aquí estados diferentes.

La inversión basada del lector de prefijos permite evaluar asimismo la
información de orientación que omite su cociente. Hay $190$ pares en
grado uno y $324$ en grado dos; cada par equiponderado aporta $\log2$
por su masa. Por tanto,
\[
 I_{\rm or}(q)=\frac{380q+648q^2}{Z_N(q)}\log2.
\]
La restitución de la etiqueta de orientación recupera esa información.
El cociente y su lectura entrópica quedan distinguidos de una operación
física de borrado.

\subsection{Dilatación energética positiva de la memoria}
\label{sec:ii-ter27-memoria}

Como $N,L$ conmutan, $Q_c=N-L$ tiene valores $0,-1,2,1,0$ en las
celdas anteriores. Sea $J_m=\operatorname{diag}(0,1,2,3)$ y $g>0$.
Defínase la isometría
\begin{equation}
 V|n,l,i\rangle=|n,l,i\rangle\otimes|n-l+1\rangle,
 \quad H_F=gL,\quad H_m=gJ_m.
 \label{eq:ii-ter27-lift}
\end{equation}
Los niveles auxiliares son $1,0,3,2,1$. La primera coordenada conserva
la base completa, por lo que $V^*V=I$. En cada vector,
\begin{equation}
 (H_F\otimes I+I\otimes H_m)V=Vg(N+I),
 \label{eq:ii-ter27-energia-memoria}
\end{equation}
porque $l+(n-l+1)=n+1$. La energía auxiliar es no negativa y las
energías totales son $g,2g,3g,3g,3g$. La imagen es reductora para este
Hamiltoniano diagonal. Para $q=e^{-\beta g}$, el cálculo funcional da
\[
 V^*e^{-\beta H_{\rm tot}}V=q^{N+I},\qquad
 \frac{V^*e^{-\beta H_{\rm tot}}V}
 {\operatorname{Tr}(V^*e^{-\beta H_{\rm tot}}V)}=\frac{q^N}{Z_N(q)}.
\]
El filtro $M_q=q^{(Q_c+I)/2}$ satisface
$V^*(I\otimes q^{J_m})V=M_q^*M_q$ y
$q^L M_q^*M_q=q^{N+I}$. Para $0<q<1$, es contractivo y realiza el
cambio condicionado de $\rho_L=q^L/Z_L$ a $\rho_N=q^N/Z_N$ con
probabilidad $qZ_N/Z_L$. Registrar su resultado conserva también la
rama complementaria. La isometría $V$ conserva las coherencias; la traza
parcial sobre la memoria constituye otra operación que elimina sólo las
coherencias entre distintos valores de $Q_c$.
El desplazamiento $+1$ es el menor entero que hace no negativos todos
los niveles auxiliares. Un desplazamiento común mayor conserva el estado
normalizado y cambia el origen energético registrado.

\subsection{Lazo de memoria y espectro decimal}
\label{sec:ii-ter27-espectro}

La realización gaussiana recibe la descomposición colectiva $8:1$ de
las nueve copias y el lazo unitario de memoria. Ambos modos se preparan
en la referencia gaussiana pura centrada de covarianza conjunta $I\oplus I$,
en las mismas coordenadas simplécticas en que se escriben las matrices
siguientes:
\[
 Q_8=\frac13\begin{pmatrix}\sqrt8 I&-I\\ I&\sqrt8 I\end{pmatrix},
 \quad H_1=\begin{pmatrix}2&1\\1&1\end{pmatrix},
 \quad U_{H_1}=Q_8^T\operatorname{diag}(I,H_1)Q_8.
\]
Estas matrices reciben su preparación y su lazo del transporte de
sistema y memoria; el cálculo térmico se efectúa después. El bloque
visible de $U_{H_1}U_{H_1}^T$ es
\[
 W=\frac{8I+H_1H_1^T}{9}
 =\frac19\begin{pmatrix}13&3\\3&10\end{pmatrix},
 \qquad \det W=\frac{121}{81},\quad \nu=\frac{11}{9}.
\]
En dimensión dos, $S=\sqrt\nu W^{-1/2}$ tiene determinante uno y es
simpléctica, con $SWS^T=\nu I$. Su realización metapléctica reduce el
estado gaussiano a esa covarianza isotrópica. El estado geométrico
$(1-r)\sum_{j\geq0}r^j|j\rangle\langle j|$ tiene ocupación media
$r/(1-r)$ y covarianza $(1+r)/(1-r)I$. La unicidad del estado gaussiano
centrado por su covarianza implica
\begin{equation}
 r=\frac{\nu-1}{\nu+1}=\frac1{10},\qquad
 \rho_r=(1-r)\sum_{j\geq0}r^j|j\rangle\langle j|.
 \label{eq:ii-ter27-williamson}
\end{equation}
La razón decimal resulta del lazo y de la preparación indicados. Por
ejemplo, $H_0=I$ produciría $\nu=1$ y $r=0$; el resultado conserva
así su dependencia efectiva del operador de memoria.

La división euclídea $j=3w+b$, $0\leq b<3$, define un unitario
$U_3|j\rangle=|w,b\rangle$ con inversa $|w,b\rangle\mapsto|3w+b\rangle$.
La factorización $1-r^3=(1-r)(1+r+r^2)$ demuestra término a término
\begin{equation}
 U_3\rho_rU_3^*=\rho_q\otimes\sigma_r,
 \quad q=r^3=\frac1{1000},\quad
 \sigma_r=\frac{\operatorname{diag}(100,10,1)}{111}.
 \label{eq:ii-ter27-agrupamiento}
\end{equation}
El desplazamiento unilateral se transporta a
\[
 T=I\otimes(|1\rangle\langle0|+|2\rangle\langle1|)
       +S_w\otimes|0\rangle\langle2|,\qquad T^3=S_w\otimes I_3.
\]
La aplicación a cada vector de base prueba esa identidad y muestra el
acarreo del residuo al cociente. Se cumple $T^*T=I$ y
$TT^*=I-|0,0\rangle\langle0,0|$; el retorno del residuo conserva el
avance de ocupación. El mapa unitario transporta también toda purificación
al actuar como $U_3\otimes I$ sobre el sistema ampliado.

Para la sección energética $H_W=\epsilon_a(j+1/2)$, el transporte da
\[
 U_3H_WU_3^*=\epsilon_a(3N_w+B+\tfrac12 I),
 \qquad B=\operatorname{diag}(0,1,2).
\]
La identidad diagonal conserva el dominio autoadjunto determinado por
la sumabilidad cuadrática de la energía. A $\beta_W\epsilon_a=\log10$,
el cociente tiene espaciamiento $3\epsilon_a$ y razón $q$, conservando
el mismo termómetro. Sus medias satisfacen
$\langle N_w\rangle=1/999$, $\langle B\rangle=4/37$ y
$3\langle N_w\rangle+\langle B\rangle=1/9$.
Además $s(\rho_r)=s(\rho_q)+s(\sigma_r)$, por la factorización exacta.

Esta precisión fija $g=3\epsilon_a$ en la composición de
\eqref{eq:ii-ter27-lift} con el cociente de Williamson. Si se prepara
$\rho_L\otimes\rho_q$ y se registra la proyección a la imagen de $V$,
\[
 V^*(\rho_L\otimes\rho_q)V=\frac{1-q}{Z_L}q^{N+I}.
\]
La probabilidad es $(1-q)qZ_N/Z_L$ y el estado condicionado es $\rho_N$.
El factor $\sigma_r$ se conserva mediante $V\otimes I_3$; la energía
total transportada es
$3\epsilon_a(N+1)+\epsilon_a(B+1/2)$.

\subsection{Sección espectral marcada e intensidad completa}
\label{sec:ii-ter27-marcador}

La composición con la capacidad utiliza el mapa declarado
$C_{s_n}\mapsto j=C_{s_n}$. Por \eqref{eq:ii-ter27-localizacion},
los niveles son $23+3n=3(7+n)+2$. La sección exige conjuntamente
residuo $2$ y cociente al menos $7$; el residuo aislado admitiría también
$2,5,8,\ldots$. Su proyector y probabilidad son
\[
 P_{23}=\sum_{n\geq0}|23+3n\rangle\langle23+3n|,
 \qquad p_{23}=\frac{(1-r)r^{23}}{1-r^3}.
\]
Condicionar y trasladar $|23+3n\rangle$ a $|n\rangle$ produce $\rho_q$.
El peso $p_{23}$ registra el umbral que desaparece de la densidad normalizada.

Para incorporar las multiplicidades $d=(1,380,649)$, sea
$\iota|n,a\rangle=|23+3n\rangle\otimes|a\rangle$ en una memoria
de dimensión $649$, con $1\leq a\leq d_n$. Es una isometría.
Comprimir $\rho_r\otimes I_{649}/649$ a su imagen da masa
\[
 p_{\mathcal D}=\frac{1-r}{649}r^{23}Z_N(q)
\]
y estado condicionado $\rho_N(q)$. Si $\lambda_j=(1-r)r^j$,
la intensidad relativa a la población de referencia $\lambda_0$ es
\begin{equation}
 B_*:=\frac{\lambda_{23}+380\lambda_{26}+649\lambda_{29}}{\lambda_0}
 =10^{-23}\left(1+\frac{380}{1000}+\frac{649}{1000^2}\right)
 =1.380649\times10^{-23}.
 \label{eq:ii-ter27-intensidad}
\end{equation}
La fórmula conserva espectro, umbral, multiplicidades y referencia.
La fórmula histórica era conocida durante la construcción de este lector;
el desarrollo operatorio precisa su realización y sus condiciones, en
lugar de atribuirle retrospectivamente una predicción independiente.
La interpretación como aplicación entre las rectas de energía y temperatura
requiere componer las lecturas dimensionales, objeto del desarrollo siguiente.

\subsection{Corrientes de frontera, compresión y rigidez del desfase}
\label{sec:ii-ter27-corrientes}

La incidencia de área ya desarrollada en \ref{sec:ii-area-incidencia-explicita}
proporciona $18+18$ enlaces tríticos. Sean $u=N_{90}$, $v=N_{120}$,
$a=u+v$ y
\[
 F=3u+4v,\qquad H_R=\epsilon_0 F,
 \quad\epsilon_0=30\Delta\tau_R>0,
 \quad\tau_R=\frac{\hbar c}{2\pi R}.
\]
El área física es $\delta a\,a$, con $\delta a$ fijada por la incidencia
y Barbero. La procedencia permanece en $u,v$; la misma área puede tener
energías diferentes. Las siguientes corrientes utilizan esa estructura
previa sin redefinir su Hamiltoniano.

En un enlace, $T^+=|1\rangle\langle0|+|2\rangle\langle1|$ y
$T^-=(T^+)^*$ satisfacen $[n,T^\pm]=\pm T^\pm$. Defínanse
\[
 J_{ef}=T^+_{120,f}T^-_{90,e},\qquad
 B_c=\left(\sum_eT^+_{90,e}\right)^3
     \left(\sum_fT^-_{120,f}\right)^2.
\]
La regla $[F,AB]=[F,A]B+A[F,B]$ prueba
\[
 [F,J_{ef}]=J_{ef},\quad[a,J_{ef}]=0,
 \qquad[F,B_c]=B_c,\quad[a,B_c]=B_c.
\]
La secuencia de ocupaciones $(0,2)\to(3,0)\to(2,1)$ tiene energías
primitivas $8,9,10$ y áreas $2,3,3$. El mismo salto energético conserva
dos acciones areales distinguibles.

Los coeficientes $c_j=[z^j](1+z+z^2)^{18}$ cuentan las ocupaciones de
cada canal: $c_0=1,c_1=18,c_2=171,c_3=1122$. El carácter de frontera
\[
 P(x)=(1+x^3+x^6)^{18}(1+x^4+x^8)^{18}
\]
da multiplicidades $g_7=324,g_8=171,g_9=1122,g_{10}=3078$.
Para los rangos $d=(1,380,649)$, el primer triple admisible de energías
consecutivas comienza en $8$. Los inicios menores se excluyen al usar
$g_0=1$, $g_1=g_2=g_5=0$, $g_3=g_4=18$, $g_6=171$, $g_7=324$ y
$g_8=171$: en cada triple falla al menos uno de los rangos requeridos.

Una compresión conectada se construye escogiendo un estado $(0,2)$,
$380$ estados $(3,0)$ y $649$ vecinos distintos de éstos por las $J_{ef}$.
La existencia de suficientes vecinos tiene una prueba de conteo: cada
ocupación de suma tres posee al menos dos predecesores de suma dos, y
cada predecesor posee como máximo $18$ sucesores. Los $380$ estados
producen, por tanto, al menos $\lceil760/18\rceil=43$ predecesores
distintos y $18\cdot43=774$ vecinos $(2,1)$. Se pueden escoger $649$.
La corriente $B_c$ conecta el estado inicial con todos los $380$ estados
intermedios; cada vecino final tiene un antecedente seleccionado. El grafo
es conectado. Sus etiquetas y la selección de subespacio se conservan.
La isometría basada $\iota_N$ satisface
\begin{equation}
 H_R\iota_N=\iota_N\epsilon_0(8I+N).
 \label{eq:ii-ter27-compresion}
\end{equation}
La prueba consiste en evaluar ambos lados en cada vector de base.

Sean $P_F=\mathbf1_{N=1}$, $B_N:\mathcal D_0\to\mathcal D_1$ y
$J_N:\mathcal D_1\to\mathcal D_2$ las corrientes comprimidas no nulas.
Para $H_\delta=\epsilon_0(8I+N)+\delta P_F$,
\[
 [H_\delta,B_N]=(\epsilon_0+\delta)B_N,
 \qquad[H_\delta,J_N]=(\epsilon_0-\delta)J_N.
\]
Conservar ambos saltos, o exigir su igualdad, fuerza $\delta=0$.
Más generalmente, una corrección diagonal conserva cada salto si sus
valores son iguales en los extremos de cada arista. La conectividad la
hace constante. Esta rigidez clasifica correcciones diagonales; su prueba
no atribuye esa clasificación a todos los operadores no diagonales.

Para $x=e^{-\epsilon_0/\tau}$ y $\tau>0$, la compresión de Gibbs es
$\iota_N^*e^{-H_R/\tau}\iota_N=x^8x^N$. Su probabilidad de selección
desde el borde es $x^8Z_N(x)/P(x)$. Registrar selección y complemento
mediante $P_N\otimes|s\rangle+(I-P_N)\otimes|f\rangle$ es isométrico
y conserva la energía si las banderas son degeneradas. En cada arista
ascendente, tasas $xr_e$ y tasa inversa $r_e>0$ satisfacen balance
detallado para los pesos $x^N$. La conectividad prueba la unicidad
estacionaria; las masas de los grados son $d_nx^n/Z_N(x)$. Se ha
especificado así una preparación térmica con baño, no una evolución
unitaria aislada.

\Needspace{29\baselineskip}
\subsection{Equivalencia térmica con referencia y escala transportadas}
\label{sec:ii-ter27-equivalencia}

\begin{theorem}[Transporte del exponente térmico centrado]
\label{thm:ii-ter27-transporte}
Sean $H_1,H_2$ autoadjuntos finitos, $\tau_i>0$ y $T$ una isometría
con imagen reductora para $H_2$. Las trazas en el segundo espacio se
toman sobre esa imagen. Sea $\sigma_{\mathrm{ref}}$ un estado normalizado
del primer espacio, y sean las energías de referencia transportadas
\[
 e_1=\operatorname{Tr}(\sigma_{\mathrm{ref}}H_1),\qquad
 e_2=\operatorname{Tr}(T\sigma_{\mathrm{ref}}T^*H_2).
\]
Si
\[
 \frac{H_2-e_2I}{\tau_2}T
 =T\frac{H_1-e_1I}{\tau_1},
\]
se conservan la partición centrada, el estado de Gibbs, la entropía y
las probabilidades de todos los observables transportados, incluido el
marcador. Para estados fieles, la igualdad de estados transportados
implica recíprocamente esa identidad con las referencias así definidas.
\end{theorem}
\begin{proof}
El cálculo funcional entrelaza las exponenciales de ambos exponentes.
La invariancia de la traza por isometría sobre la imagen da igualdad de
las particiones centradas; al normalizar resulta $\rho_2=T\rho_1T^*$.
La ciclicidad de la traza transporta observables y el cálculo espectral
transporta la entropía. Recíprocamente, tomar logaritmos de los estados
fieles produce igualdad de $H_i/\tau_i$ salvo un escalar, igual a la
diferencia de sus logaritmos de partición. Tomar la esperanza en
$\sigma_{\mathrm{ref}}$ identifica ese escalar con
$e_2/\tau_2-e_1/\tau_1$; restarlo da la identidad centrada enunciada.
\end{proof}

En el lector de capacidad, $C=23I+3N$ y
\begin{equation}
 H_{\rm cap}=\frac{\epsilon_a\log10}{\log3}C,
 \quad\tau_{\rm clk}=\frac{\epsilon_a}{\log3},\qquad
 H_{W,\rm rel}=\epsilon_a C,
 \quad\tau_W=\frac{\epsilon_a}{\log10}.
 \label{eq:ii-ter27-hcap}
\end{equation}
Ambos exponentes son $C\log10$, aunque energías y temperaturas tengan
escalas diferentes. En la compresión del borde, la preparación
$\tau_*=\epsilon_0/\log1000$ satisface
\[
 \frac{H_R-8\epsilon_0I}{\tau_*}\iota_N
 =\iota_N\frac{H_{\rm cap}-23\tau_{\rm clk}\log10\,I}{\tau_{\rm clk}}
 =\iota_N N\log1000.
\]
La actividad nativa de horizonte sigue siendo $e^{-30\Delta}$; esta
preparación tiene razón $\tau_*/\tau_R=30\Delta/\log1000$. Al variar
$R$ cambian $\epsilon_0$ y $\tau_R$ conjuntamente, conservando su cociente.
Sin centrar, los exponentes de borde y capacidad difieren en
$\log10\,I$, pues $8\log1000-23\log10=\log10$. Sus trazas son
$10^{-24}Z_N$ y $10^{-23}Z_N$. El factor $1/10$ también aparece en
sus referencias y debe conservarse al comparar intensidades absolutas.

Si $\iota_1,\iota_2$ realizan los mismos rangos, grados y marcador,
$U=\iota_2\iota_1^*$ es unitario entre sus imágenes y entrelaza
Hamiltoniano y marcador. Por el teorema anterior, cualquier elección de
base admisible da el mismo lector, todos sus momentos espectrales y sus
derivadas térmicas. La unicidad del representante microscópico es
innecesaria para esa conclusión. Comparar además área, corrientes o
tiempos de relajación exige transportar esos observables adicionales.

Finalmente, para $P_0=|\Omega_{\rm ext}\rangle\langle\Omega_{\rm ext}|$
y $\rho_N=q^N/Z_N$, se tiene
\[
 p_0=\operatorname{Tr}(P_0\rho_N)=Z_N^{-1},\qquad
 Z_N=1+\frac{\operatorname{Var}_{\rho_N}(P_0)}{p_0^2}.
\]
La segunda identidad usa $P_0^2=P_0$, de modo que la varianza es
$p_0(1-p_0)$. El estado normalizado y el marcador recuperan la traza
relativa; el umbral registrado recupera el factor de
\eqref{eq:ii-ter27-intensidad}. El resultado conserva así la estructura,
la preparación, la referencia y la memoria necesarias para su composición
con la lectura de acción y temperatura.
